\documentclass[12pt]{article}
\usepackage[utf8]{inputenc}
\usepackage{graphicx}
\usepackage{amsmath}
\usepackage{amssymb}
\usepackage{mathrsfs}
\usepackage{pgfornament}
\usepackage[many]{tcolorbox}
\usepackage[margin = 0.5in]{geometry}
\usepackage{collectbox}
\usepackage{mdframed}
\usepackage{xcolor}
\usepackage{mathtools}
\DeclarePairedDelimiter\ceil{\lceil}{\rceil}
\DeclarePairedDelimiter\floor{\lfloor}{\rfloor}
\newcommand\norm[1]{\left\lVert#1\right\rVert}
\newcommand\textbox[1]{%
  \parbox{.333\textwidth}{#1}%
}

\linespread{2}

\makeatletter
\newcommand{\mybox}{%
    \collectbox{%
        \setlength{\fboxsep}{2pt}%
        \fbox{\BOXCONTENT}%
    }%
}
\makeatother

\newcommand{\innerprod}[2]{\left<#1, #2\right>}

\usepackage{listings}
\usepackage{color}

\definecolor{dkgreen}{rgb}{0,0.6,0}
\definecolor{gray}{rgb}{0.5,0.5,0.5}
\definecolor{mauve}{rgb}{0.58,0,0.82}

\lstset{frame=tb,
  language=R,
  aboveskip=3mm,
  belowskip=3mm,
  showstringspaces=false,
  columns=flexible,
  basicstyle={\small\ttfamily},
  numbers=none,
  numberstyle=\tiny\color{gray},
  keywordstyle=\color{blue},
  commentstyle=\color{dkgreen},
  stringstyle=\color{mauve},
  breaklines=true,
  breakatwhitespace=true,
  tabsize=3
}
\title{MAT 523 HW 10}
\begin{document}

\hfill April 17, 2019
\begin{center}
\begin{large}
\textbf{MAT 524: Functions of a Real Variable II \\
\vspace{-5mm}
Homework X \\}
\end{large}
\vspace{-2mm}
Nolan R. H. Gagnon \\
\end{center}
\vspace{-0.5cm}
\noindent \textbf{{\huge [}{\Large [}} Problem One \textbf{{\Large ]}{\huge ]}}{\Large \textbf{:}}

\vspace{3mm}

\noindent\mybox{\colorbox{yellow}{\textbf{Problem Statement}}}\hspace{3mm}\hrulefill

\noindent For $f,g\in L^{1}(\mathbb{R})$, define their convolution by $$f * g(x) = \int_{\mathbb{R}} f(x-t)g(t) dt.$$ It can be shown that $f*g(x)$ exists for almost all $x$, and that $f*g\in L^{1}(\mathbb{R})$. Granting this, prove that \textbf{(a)} $f*g = g*f$ and \textbf{(b)} $\widehat{f *g} = \hat{f}\hat{g}.$

\noindent\mybox{\colorbox{yellow}{\textbf{Solution}}}\hspace{3mm}\hrulefill

\noindent \textbf{(a)} We have
\begin{align}
    f * g(x) &= \int_\mathbb{R} f(x-t)g(t)dt. \label{eq1.1}
\end{align}
Let $u=x-t$. Then, if we consider $x$ to be fixed, we can write $du = -dt$. Furthermore, we have $t=x-u$. With this substitution, we have
\begin{align}
    \int_\mathbb{R} f(x-t)g(t)dt &= \int_{\infty}^{-\infty} f(u)g(x-u)(-du) \\
    &=\int_{-\infty}^{\infty}g(x-u)f(u)du \\
    &=\int_\mathbb{R}g(x-u)f(u)du \\
    &= g * f(x).
\end{align}

\vspace{5mm}

\noindent\textbf{(b)} First note that since $f*g\in L^1(\mathbb{R})$, it follows that $\int_{\mathbb{R}} |f*g| < \infty$, i.e., the integral converges absolutely. The Fourier transform of $f*g$ is
\begin{align}
    \widehat{f*g}(y) &= \int_{\mathbb{R}}\left(\int_{\mathbb{R}}f(x-t)g(t)dt\right)e^{-2\pi ixy}dx \\
    &= \int_\mathbb{R}\int_\mathbb{R} f(x-t)g(t)e^{-2\pi ixy}dtdx.
\end{align}
Let $u=x-t$. Then $x=u+t$ and $du=dx$. Then we have, by Fubini's Theorem,
\begin{align}
    \int_\mathbb{R}\int_\mathbb{R} f(x-t)g(t)e^{-2\pi ixy}dtdx &= \int_\mathbb{R}\int_\mathbb{R}f(u)g(t)e^{-2\pi(u+t)y}dtdu \\
    &= \left(\int_\mathbb{R}f(u)e^{-2\pi i uy}du\right)\left(\int_\mathbb{R}g(t)e^{-2\pi ity}dt\right) \\
    &= \hat{f}(y)\hat{g}(y).
\end{align}

\hfill$\blacksquare$

\vfill

\begin{center}
\pgfornament[width = 50mm, color = black, symmetry = h]{83}\\
\vspace{2mm}
\vspace{-0.65cm}
\hrulefill \\
\vspace{-0.95cm}
\hrulefill
\end{center}

\end{document}
